\paragraph{\trMovingBoundaryCondition。}\label{sub:moving-boundary-conditions}
设 $S$ 是带\trStriplikeEnd{}的\trRiemann\trSurface。$S$ 的一组\trMovingBoundaryCondition{}由一族%
\trExact\trLagrange\trSubmanifold{}$L_z\subset M$ 组成，它们\trSmooth{}地依赖于 $z\in\partial S$，但在
$z\in\partial S\cap\epsilon_\zeta(Z^\pm)$ 上为\trLocallyConstant，于是每个无穷远点仍有一对与之
相关联的\trLagrange\trSubmanifold{}$(L_{\zeta,0},L_{\zeta,1})$。在\trPerturbationDatum{}$(K,J)$ 的
定义中，\eqref{eq:k-boundary-condition}应代之以下面这个稍更复杂的条件：
\begin{equation}\label{eq:moving-boundary-condition}
d(K(\xi)|L_z)=\phi(\cdot,\partial_\xi L_z)\in\Omega^1(L_z)
\qquad\text{对任一 }\xi\in T(\partial S)_z.
\end{equation}
这里 $\partial_\xi L_z$ 是 $C^\infty(TM|L_z)$ 的、切于该\trLagrange\trSubmanifold{}之无穷小%
\trDeformation{}的任一\trSection，于是\eqref{eq:moving-boundary-condition} 的右端就是与
这样一种\trDeformation{}相对应的标准（\trExact）1-形式。接下来便可以同从前一样
考虑同一个方程\eqref{eq:inhomogeneous-equation}。对 $K$ 的这个条件使人能把%
\eqref{eq:perturbed-action} 中的边界项积出去，而理论的其余部分不须任何实质性
调整即可通行，从而得到广义化的映射\eqref{eq:tqft-operation}，它们满足 TQFT
公理的显然变体。事实上，这种一般性的增加远比初看起来要小。为看出这一点，
假设 $\partial S$ 的一切\trConnectedComponent{}都是开区间。那么存在一族\trAutomorphism
$\phi_z\in\mathrm{Aut}(M,\partial M)$，由 $z\in S$ \trParametrization，使得 $\phi_z(L_z)$ 为%
\trLocallyConstant，并且在诸\trStriplikeEnd{}上成立 $(\partial/\partial s)%
\phi_{\epsilon_\zeta(s,t)}=0$。坐标变换
\begin{equation}\label{eq:moving-coordinate-change}
u(z)\longmapsto\phi_z(u(z))
\end{equation}
保持\eqref{eq:inhomogeneous-equation} 这类方程，但把\trMovingBoundaryCondition{}化归
为常值的。尽管如此，用\trMovingBoundaryCondition{}来思考常常是方便的，正如下面的
应用所表明的那样。设 $\{L_s\}_{0\le s\le 1}$ 是一个\trExact\trLagrange\trIsotopy。取 $S=H$
并配以\trMovingBoundaryCondition{}$L_{\psi(s)}$，其中 $\psi\in C^\infty(\RR,[0;1])$ 满足：
当 $s\ll 0$ 时
$\psi(s)=0$、当 $s\gg 0$ 时 $\psi(s)=1$。这就给出一个特异元素
\begin{equation}\label{eq:isotopy-characteristic-element}
\Phi_H\in HF^{\mathrm{pr}}(L_0,L_1),
\end{equation}
它事实上就是\cref{sub:donaldson-fukaya-category} 中提到的那一个。那里陈述的
全部性质都可由标准的\trHomotopy{}与\trGluing{}论证得到。例如，把%
\eqref{eq:isotopy-characteristic-element} 与\trTriangleProduct{}作\trComposition，则对任意的
$L$ 得到\trIsomorphism{}$HF^{\mathrm{pr}}(L,L_0)\to HF^{\mathrm{pr}}(L,L_1)$ 与
$HF^{\mathrm{pr}}(L_1,L)\to HF^{\mathrm{pr}}(L_0,L)$。事实上，细察这个\trGluing{}论证
并对之施加一个变换\eqref{eq:moving-coordinate-change}，就能把它们化归为先前
考虑过的\trContinuationMap。
